


\begin{center}
    \textbf{Disclaimer:} 
    This document was written by Markus Renoldner, who is responsible for any mistakes.
\end{center}

\vspace{0.5cm}

\tableofcontents

\newpage
% \listoftodos 

% \vspace{0.5cm}
\newpage
\section{Introduction}

These are the lecture notes based on the lecture MATH-344 about Riemannian geometry given by Prof. Giorgios Moschidis at EPFL during the spring semester of 2026. The lecture covers the following topics:
\begin{itemize}
    \item Riemannian metrics: Inner products on tangent spaces enabling length and distance measurements
    \item Connections and geodesics: Generalizing straight lines to curved manifolds
    \item Curvature tensors: Quantifying how manifolds bend
    \item Metric structure: Including the Hopf-Rinow theorem relating geometry to topology
    \item Hypersurfaces: Geometry of submanifolds
    \item Comparison theorems: Relating geometry to topology
\end{itemize}

The course extends Euclidean geometry concepts to the broader setting of Riemannian manifolds, essential for applications in PDEs, differential geometry, and theoretical physics. As a reference, the reader may consult \cite{doCarmo1992RiemannianGeometry}. The prerequisite is differential topology at the level of the Differential Geometry II course at EPFL, \cite{tsakanikas2024dg2}, which is based on \cite{lee2003smooth}. Another popular reference is \cite{lee2018riemannian}.











